\section{ MC simulation of a Lennard-Jones fluid in the NVT ensemble}

\subsection{}
we know that the inter molecular Lennard-Jonnes potential is described by

\begin{equation}
	u(r_{ij}) =4\epsilon\left[ \left(\frac{\sigma}{r_{ij}}\right)^{12}-\left(\frac{\sigma}{r_{ij}}\right)^6\right] 
\end{equation}

now for simulation purposes a cutoff point is defined, on this point the potential is almost 0

\begin{equation}
		u_c(r_{ij}) = 
		\begin{cases}
			u(r_{ij})-e_c\ & \text{if}\ r_{ij} \le r_c\\
			0\ & \text{if}\  r_{ij} > r_c
		\end{cases}
\end{equation}

taking the derivative to get the force we get

\begin{equation}
	\begin{split}
		f(r_{ij}) & = - \nabla u(r_{ij})\\
		& =  \frac{48 \epsilon}{r_{ij}} \left[ \left(\ \frac{\sigma}{r_{ij}}\right)^{12}-\frac{1}{2}\left(\frac{\sigma}{r_{ij}}\right)^6\right] 
	\end{split}
\end{equation}

Meaning that the force function used for simulating can be described by
\begin{equation}
	f_c(r_{ij}) = 
	\begin{cases}
		f(r_{ij})\ & \text{if}\ r_{ij} \le r_c\\
		0\ & \text{if}\  r_{ij} > r_c
	\end{cases}
\end{equation}



We know that the equation for the pressure is
\begin{equation}
	\left\langle P \right\rangle= \rho k_B T + \frac{1}{3V} \left\langle \sum_{i<j} \vec{f(r_{ij})} \cdot \vec{r_{ij}} \right\rangle
\end{equation}


Using all these equations we can get the pressure



\subsection{}

We also want to see what the chemical potential is, I.E. how much energy does it cost to add one particle into the system. 

\begin{equation}
	\mu_{ex} = -k_B T \ln \left[ \sum_{i=1}^{N_\text{test}} \frac{\exp \left\{-\beta \Delta U_i^+\right\}}{N_\text{test}}\right]
\end{equation}


